\chapter{Introducción}
\label{sec:man-intro}

\section{Qué es GraphHunter}

GraphHunter es una herramienta de caza de amenazas sobre grafos
temporales. Toma logs de seguridad de distintas fuentes (Sysmon,
Sentinel, FortiAnalyzer, IIS, Cognito, CSVs genéricos) y construye
un grafo donde los nodos son entidades (usuarios, hosts, IPs,
procesos, archivos, dominios) y las aristas son relaciones
observadas en el tiempo (autenticar, ejecutar, conectar, escribir).

Sobre ese grafo, el analista describe \emph{hipótesis} con un DSL
compacto y el sistema devuelve los paths que las satisfacen,
scoreados por anomalía.

\textbf{Véase también} el libro \emph{Diseño completo}
(\filepath{docs/diseno/diseno.pdf}) si querés entender qué pasa por
detrás cuando hacés clic en ``Run hunt'': cómo el motor elige el
operador de matching, qué índices consulta, por qué el primer
\emph{hunt} sobre una sesión recién cargada tarda más, y qué nuevas
\emph{shapes} de hipótesis acepta el DSL.

\section{Cuándo usarlo}

\begin{itemize}
    \item \textbf{Caso típico}: un incidente (alerta, ticket CSIRT,
    sospecha de lateral movement) que requiere reconstruir qué pasó
    y quién participó a partir de logs dispersos.
    \item \textbf{Investigación proactiva}: formular hipótesis de
    TTPs ATT\&CK sobre logs históricos.
    \item \textbf{Análisis post-mortem}: documentar una cadena de
    eventos y exportar IOCs u OCSF para entregar al cliente.
\end{itemize}

GraphHunter no reemplaza al SIEM. El SIEM es bueno para búsquedas
textuales y alertas por regla; GraphHunter es bueno para
preguntas estructurales: ``¿qué usuarios que autenticaron anoche
también escribieron en rutas sensibles dentro de la siguiente
hora?''.

\section{Flujo típico de un turno}

\begin{enumerate}
    \item Abrir la app, crear una sesión nueva.
    \item Ingestar los logs relevantes (drag-and-drop, o conectar a
    Sentinel).
    \item Inspeccionar el catálogo: ¿qué tipos de entidad apareció?
    ¿qué relaciones se observaron?
    \item Escribir una hipótesis en DSL y correr el hunt.
    \item Revisar los paths devueltos ordenados por score;
    explorar los de mayor rareza en el \emph{modo explorador}.
    \item Etiquetar entidades relevantes, agregar notas y exportar
    resultados.
\end{enumerate}

El capítulo~\ref{sec:man-quickstart} te guía por este flujo
end-to-end con datos demo en unos 10 minutos.

\section{Este manual}

Este libro es para el analista. No requiere saber Rust, TypeScript
ni los internals. Si querés entender cómo está construido
GraphHunter, ver el libro de Arquitectura. Si necesitás la
referencia de API o DSL con detalle de contratos, ver
Especificación.

\section{Convenciones}

\begin{itemize}
    \item \code{Texto monoespaciado} refiere a comandos, endpoints,
    código o valores literales.
    \item \filepath{rutas/entre/barras} son rutas del proyecto.
    \item Los ejemplos asumen que la app corre en
    \code{127.0.0.1:37891}. Si configuraste otro puerto, ajustá.
\end{itemize}
